\documentclass[]{article}

\usepackage{amsmath}
\usepackage{multirow}
\usepackage{natbib}
\usepackage{graphicx} 


\begin{document}

The formation of the first globular clusters (GCs) is hypothesized to be regulated by the atomic cooling threshold, which predicts their assembly in dark matter halos with virial temperatures exceeding $10^4$ K. We test this framework across cosmic time by comparing two distinct GC populations: the 19 clusters of the GEMS system observed at $z=9.625$ and the 5 clusters of the Sparkler system at $z=1.4$. By calculating the formation redshift ($z_{form}$) for each cluster from its published age, we map their empirical formation epochs onto the theoretical GC formation rate predicted by the model. We find the Sparkler GCs, with $z_{form}$ between 2.2 and 3.5, align with the predicted peak of formation activity, while the GEMS GCs, with $z_{form}$ between 9.7 and 19.1, populate the high-redshift tail of the same distribution, a result consistent with an observational selection effect. Furthermore, the GEMS clusters are unexpectedly more metal-rich than their lower-redshift Sparkler counterparts, implying their formation occurred within a massive and rapidly enriching host environment at cosmic dawn. The alignment of these two disparate populations with different epochs of a single theoretical framework suggests the atomic cooling threshold acts as a primary regulator of GC formation.

\end{document}
                